\section{Lorentz Invariance on Every Stratum}
\paragraph*{Theorem 2.1 (Stratified Lorentz invariance). The Minkowski norm |qn|2 = $\tau$ 2}
n is invariant under the Lorentz boost J : z-7$\to$-z-on every stratum Sn. Moreover, every change of interpretation of C+ (chirality, Weyl chamber, upper half-plane, or any other frame-dependent labeling) leaves $\tau$n and the compensator C(f) = 1 2(f + Jf) invariant.
\paragraph*{Proof. Norm invariance. Under J : qn 7$\to$bar qn (quaternion conjugate): (it)*=}
it, z* i = zi for the spectral coordinates. The boost acts as z-7$\to$-z-on the contravariant component, leaving |z-|2 unchanged. Hence |qn|2 = $\tau$ 2 n is invariant. Frame independence of C. The compensator C(f) = 1 2(f + Jf) requires only J2 = id (which holds as J is an involution). Any relabeling of C+ corresponds to a unitary transformation U on H+, but C = 1 2(id + J) is defined intrinsically from J, hence UCU -1 = C for any U commuting with J --- which includes all frame changes generated by the Sphaera-Cascade.
\paragraph*{Remark (Why the interpretation is irrelevant). The biquaternion qn can be}
read as a spinor (Clifford algebra), a Dirac bispinor (QFT), an element of the upper half-space (hyperbolic geometry), or a point in projective twistor space. In every reading, $\tau$n is an invariant and C projects onto K3. This is the precise mathematical content of the observation that "it is egal in every frame": the frame determines the interpretation, but not the structure.

